\selectlanguage{ngerman}%
\frenchspacing%
\begin{ledgroupsized}[r]{120mm}%
\footnotesize%
\pstart%
\noindent\textbf{Überlieferung:}%
\pend%
\end{ledgroupsized}%
\begin{ledgroupsized}[r]{114mm}%
\footnotesize%
\pstart%
\parindent -6mm%
\makebox[6mm][l]{\textit{L}}%
Auszüge mit Bemerkungen aus
\textsc{C.\:F.\:M.\;Dechales},%
\indexN{\textls{Dechales} (Chalesius), Claude François Milliet 1621\textendash1678}
\textit{Cursus seu mundus mathematicus} (3~Bde, Lyon 1674)%
:\cite{00124}
LH~XXXV\;9,\;17 Bl.\;5--6.
Ein Bogen 2\textsuperscript{o}.
Eine Spalte auf Bl.\;6\,r\textsuperscript{o}.
Bl.\;5\,r\textsuperscript{o} bis Bl.\;6\,r\textsuperscript{o} überliefern N.\;\ref{41182-1} im vorliegenden Band (Auszüge mit Bemerkungen zur Pneumatik aus demselben Werk).
Bl.\;6\,v\textsuperscript{o} ist unbeschrieben.
\textemdash\
Weitere Auszüge aus Dechales' \textit{Cursus} überliefern im Leibniz-Nachlass die Manuskripte:
LH~XXXV\;9,\;17 Bl.\;1--4; %
LH~XXXVI Bl.\;54--57, 60--61; %
LH~XXXVII\;5 Bl.\;2--3 %
(sämtliche werden voraussichtlich in künftigen Bänden von \textit{LSB} VIII ediert).
Ein Exemplar von Dechales' \textit{Cursus} mit Gebrauchsspuren von Leibnizens Hand liegt in der GWLB Hannover unter der Signatur Nm-A\,10018:1--3 vor.
\pend%
\end{ledgroupsized}%
\vspace{5mm}%
\begin{ledgroup}%
\footnotesize%
\pstart%
\noindent%
\textbf{Datierungsgründe:}
Das vorliegende undatierte Stück N.\;\ref{41182-2} umfasst Leibnizens ausführliche Exzerpte aus \textsc{C.\:F.\:M.\;Dechales}, \textit{Cursus seu mundus mathematicus} %
tract.\;VII: \textit{Statica}, lib.\;III: \textit{De descensu gravium in planis inclinatis et funependulis} (Bd.\;I, Lyon 1674, S.\;489--524, bes. S.\;491--496).\cite{00124}
Die vorgeschlagene Datierung beruht auf der Annahme, dass die Exzerpte N.\;\ref{41182-2} zum gleichen Zeitpunkt entstanden sind wie die auf demselben Bogen überlieferten Exzerpte aus Dechales'\indexN{\textls{Dechales} (Chalesius), Claude François Milliet 1621\textendash1678} \textit{Cursus}, %
die in diesem Band unter N.\;\ref{41182-1} ediert sind.
Siehe zur Datierung von N.\;\ref{41182-1} die ausführliche Begründung auf S.\;\pageref{41182_dat}.%
\pend%
\end{ledgroup}%
\selectlanguage{latin}%
\frenchspacing%
\vspace{8mm}%
\normalsize%
\pstart%
\noindent%
\bracks{%
\href{\Scan{LH_35_09_17_006r}}%
{6\,r\textsuperscript{o}}%
}
\edtext{%
\textit{Gravia in planis inclinatis%
\indexS{planum inclinatum}%
\indexS{grave in plano inclinato}
eandem altitudinem perpendicularem habentibus%
\indexS{altitudo perpendicularis}
posita,
et ita se habentia %
ut planorum}
\edtext{%
\textit{longitudines},
sunt in aequilibrio,%
\indexS{gravia in aequilibrio}%
}{%
\lemma{\textit{longitudines},}%
\Bfootnote{%
\textit{(1)}~\textit{aequalia habent mom} %
\textit{(2)}~sunt in aequilibrio,%
~\textit{L}%
}}
seu
\edtext{%
si pondus \textit{A} ad pondus \textit{B}%
\indexS{pondus gravium in plano inclinato}
ut recta \textit{AC} ad rectam \textit{CB}.%
}{%
\lemma{pondus \textit{A} \bracks{...} rectam \textit{CB}}%
\Cfootnote{%
Vgl. das Diagramm \bracks{\textit{Fig.\:1}}.%
}}
\textit{Ingeniosissima}
\edtext{%
\textit{demonstratio Stevini},%
\indexN{\textls{Stevin} (Stevinus), Simon 1548\textendash1620}%
\indexS{demonstratio Stevini}%
}{%
\lemma{\textit{demonstratio}}%
\Cfootnote{%
Siehe \textsc{S.\;Stevin}, \textit{De beghinselen der weeghconst},
Buch~I, Teil~II, Theor.\;XI, Prop.\;XIX mit Koroll.\;I
(Leiden 1586, S.\;40\textendash42;\cite{00196}
\textit{SPW}~I, S.\;174\textendash179).\cite{02084}%
}}
alioqui posito
\edtext{%
fune gravi circumplicato \textit{ACBA},%
\indexS{funis circumplicatus gravi}%
}{%
\lemma{fune \bracks{...} \textit{ACBA}}%
\Cfootnote{%
Vgl. das Diagramm \bracks{\textit{Fig.\:2}}.%
}}
vel etiam globulis aequaliter instructo%
\indexS{funis instructus globulis}%
\bracks{,}
necessario \textit{sequetur motus perpetuus}.%
\indexS{motus perpetuus}
Idem demonstratur ex principio generali,%
\indexS{principium generale}
quia \textit{perpendicula} descensum \textit{metientia}
sunt ponderibus \textit{reciproca}.%
\indexS{perpendicula metientia descensum sunt ponderibus reciproca}
Item \textit{supposito corpus non moveri
nisi centrum gravitatis descendat}.%
\indexS{descensus centri gravitatis}%
}{%
\lemma{\textit{Gravia} \bracks{...} \textit{descendat}}%
\Cfootnote{%
Vgl. \textsc{C.\,F.\,M.\;Dechales},
\textit{Cursus seu mundus mathematicus},
tract.\:VII, lib.\;III, prop.\;IV %
(Lyon 1674, Bd.\;I, S.\;491b--492a: Zitat mit Auslassungen).%
\cite{00124}%
}}
\edtext{%
\textit{Hinc momenta gravium aequalium
in planis inclinatis consistentium},%
\indexS{planum inclinatum}%
\indexS{grave in plano inclinato}
sunt \textit{ut sinus recti inclinationum%
\indexS{sinus rectus inclinationis}%
\indexS{inclinatio plani}%
\indexS{momenta gravium in planis inclinatis sunt ut sinus inclinationum}
seu angulorum} \textit{A},\;\textit{B}.%
}{%
\lemma{\textit{Hinc} \bracks{...} \textit{angulorum} \textit{A},\;\textit{B}}%
\Cfootnote{%
Vgl. \textsc{Dechales}, \textit{Cursus}, tract.\:VII, lib.\;III, %
prop.\;VI, coroll.
(S.\;492b: Zitat mit Auslassungen).%
\cite{00124}%
}}
\pend%
\newpage%
\begin{edfigure}[0.00em][1.5em]%
\hfill%
\includegraphics[scale=0.90]{%
gesamttex/edit_VIII-4/images/LH_35_09_17_005-006_d12_006r.pdf}
\hfill%
\includegraphics[scale=0.90]{%
gesamttex/edit_VIII-4/images/LH_35_09_17_005-006_d13_006r.pdf}
\hfill%
\edcaption[0.25em]{%
\hspace*{-5mm}%
\lbrack\textit{Fig.\:1}\rbrack%
\edtext{}{%
\lemma{\lbrack\textit{Fig.\:1}\rbrack}%
\Cfootnote[nonum]{%
Vgl. a.a.O., %
prop.\;IV (S.\;491b: Detail aus dem einzigen Diagramm).%
\cite{00124}%
}%
}%
\hspace*{42mm}%
\lbrack\textit{Fig.\:2}\rbrack%
\edtext{}{%
\lemma{\lbrack\textit{Fig.\:2}\rbrack}%
\Cfootnote[nonum]{%
Vgl. a.a.O. %
(S.\;492a: erstes Diagramm).%
\cite{00124}%
}%
}%
}%
\end{edfigure}
\pstart%
\edtext{%
Sit
\edtext{%
pondus \textit{E}%
}{%
\lemma{pondus\;\textit{E}}%
\Cfootnote{%
Vgl. das Diagramm \bracks{\textit{Fig.\:3}}.%
}}
trahens \textit{D} in linea \textit{FD}%
\indexS{pondus trahens alterum in plano inclinato}%
\indexS{tractio ponderis in plano inclinato}
ipsi \textit{BC} plano inclinato parallela,%
\indexS{planum inclinatum}
et cum eo in aequilibrio,%
\indexS{pondera aequilibrata in plano inclinato}
et pondus \textit{P} trahens non parallela \textit{ID},
et tamen cum pondere \textit{D}
\edtext{%
in aequilibrio.%
\indexS{pondera aequilibrata in plano inclinato}
Ergo et pondera%
}{%
\lemma{in aequilibrio}%
\Bfootnote{%
\textit{(1)}~, debent pondera %
\textit{(2)}~, perinde est ac si pondera %
\textit{(3)}~. Ergo et pondera%
~\textit{L}%
}}
\textit{E} et \textit{P}%
\indexS{pondera aequilibrata in plano inclinato}
eorundem momentorum,%
\indexS{momentum gravis in plano inclinato}
ergo fingendo Hypomochlium \textit{G}%
\indexS{hypomochlium}
erunt reciproce ut distantiae ab eo,
\textit{GD} et
\edtext{%
\bracks{\textit{YG}},%
}{%
\lemma{\textit{YD}}%
\Bfootnote{%
\textit{L\;ändert Hrsg. nach Vorlage}%
}}
seu \textit{P} %
ad \textit{E}
ut \textit{DG} ad \textit{YG}.
\edtext{%
Jam \textit{E} ad \textit{D}
ut \textit{KG} ad \textit{DG}.%
}{%
\lemma{Jam}%
\Bfootnote{%
\textit{(1)}~\textit{D} ad \textit{E} %
\textit{(2)}~\textit{E} ad \textit{D} ut %
\textit{(a)}~\textit{DG} %
\textit{(b)}~\textit{KG} ad \textit{DG}.%
~\textit{L}%
}}
Ergo \textit{P} ad \textit{D}
ut \textit{KG} sinus inclinationis plani,%
\indexS{sinus rectus inclinationis}%
\indexS{inclinatio plani}
ad \textit{YG} sinum compl.%
\indexS{sinus complementi anguli tractionis}
\edtext{%
anguli \textit{IDF}
qui angulus tractionis est.%
\indexS{angulus tractionis}%
}{%
\lemma{anguli}%
\Bfootnote{%
\textit{(1)}~tractionis %
\textit{(2)}~\textit{IDF} qui %
\textit{(a)}~est %
\textit{(b)}~angulus tractionis est.%
~\textit{L}%
}}%
}{%
\lemma{Sit \bracks{...} est}%
\Cfootnote{%
Vgl. a.a.O., prop.\;IX %
(S.\;494a: Zusammenfassung).\cite{00124}%
}}
\pend%
\pstart%
\edtext{%
\edtext{%
\textit{BEADC} funis,%
\indexS{funis tensus}%
}{%
\lemma{\textit{BEADC} funis}%
\Cfootnote{%
Vgl. das Diagramm \bracks{\textit{Fig.\:4}}.%
}}
angulus \textit{A} rectus%
\indexS{angulus rectus}%
\bracks{,}
erunt pondera
\textit{B}, \textit{A}, \textit{C}
homologa seu proportionalia%
\indexS{pondus homologum lateri opposito}
lateribus oppositis
\textit{AD}, \textit{DE}, \textit{EA}.
Nam pondus \textit{A}
trahitur%
\indexS{pondus trahens alterum in plano inclinato}%
\indexS{tractio ponderis in plano inclinato}
\edtext{%
a pondere \textit{B},%
}{%
\lemma{a}%
\Bfootnote{%
\textit{(1)}~chorda %
\textit{(2)}~pondere \textit{B},%
~\textit{L}%
}}
quasi suspensum ex \textit{DA}%
\indexS{pondus suspensum a brachio staterae}
brachio staterae inclinatae,%
\indexS{brachium staterae inclinatae}%
\indexS{statera inclinata}
unde \textit{A} ad \textit{B}
ut \textit{AD} ad \textit{FD}
seu ut \textit{ED} ad \textit{DA}.
Ita et in altero \textit{C}.%
}{%
\lemma{\textit{BEADC} \bracks{...} altero \textit{C}}%
\Cfootnote{%
Vgl. \textsc{Dechales}, \textit{Cursus}, tract.\:VII, lib.\;III, prop.\;X %
(S.\;494b--495a: Zusammenfassung).\cite{00124}%
}}
\pend%
\begin{edfigure}[1.50em][0.00em]%
\hfill%
\includegraphics[scale=0.90]{%
gesamttex/edit_VIII-4/images/LH_35_09_17_005-006_d14_006r.pdf}
\hfill%
\includegraphics[scale=0.90]{%
gesamttex/edit_VIII-4/images/LH_35_09_17_005-006_d15_006r.pdf}%
\hfill%
\edcaption[0.5em]{%
\hspace{-1mm}%
\lbrack\textit{Fig.\:3}\rbrack%
\edtext{}{%
\lemma{\lbrack\textit{Fig.\:3}\rbrack}%
\Cfootnote[nonum]{%
Vgl. \textsc{Dechales}, \textit{Cursus}, tract.\:VII, lib.\;III, %
prop.\;IX (S.\;494a: einziges Diagramm).\cite{00124}%
}}%
\hspace{48mm}%
\lbrack\textit{Fig.\:4}\rbrack%
\edtext{}{%
\lemma{\lbrack\textit{Fig.\:4}\rbrack}%
\Cfootnote[nonum]{%
Vgl. a.a.O., prop.\;X %
(S.\;494b: einziges Diagramm).\cite{00124}%
}}%
}%
\end{edfigure}%
\newpage%
\pstart%
\edtext{%
Si
\edtext{%
angulus \textit{A}%
}{%
\lemma{angulus \textit{A} acutus}%
\Cfootnote{%
Vgl. das Diagramm \bracks{\textit{Fig.\:5}}.%
}}
\edtext{%
acutus,%
\indexS{angulus acutus}
tunc \textls{anguli tractionis}%
\indexS{angulus tractionis}
complementum erit angulus \textit{A},%
\indexS{complementum anguli tractionis}%
}{%
\lemma{acutus,}%
\Bfootnote{%
\textit{(1)}~erit \textit{A} complementum %
\textit{(2)}~tunc \textls{anguli} %
\bracks{...} angulus \textit{A},%
~\textit{L}%
}}
ergo per dicta
erit \textit{A} ad \textit{B}
ut \textit{EI} ad \textit{EG}.
Sed triangula \textit{GAN}, \textit{IEG} sunt similia.%
\indexS{triangula similia}
Nam ob angulos \textit{EGA}, \textit{EIA} rectos,%
\indexS{angulus rectus}
puncta \textit{A}, \textit{I}, \textit{G}, \textit{E}
erunt in semicirculo super \textit{AE}.
\textit{%
Cumque anguli
\edtext{%
\lbrack\textit{GAI}, \textit{GEI}\rbrack%
}{%
\lemma{\textit{GAE}, \textit{GIE}}%
\Bfootnote{%
\textit{L\;ändert Hrsg. nach Vorlage}%
}}
eidem arcui \textit{GI} insistant,%
\indexS{anguli insistentes eidem arcui}%
\indexS{arcus circuli}
erunt aequales,
item anguli \textit{GIE}, \textit{GAE}
insistentes eidem arcui \textit{GE},
et cum lineae \textit{NG}, \textit{AE}
sint parallelae
anguli alterni \textit{NGA}, \textit{GAE}
erunt aequales,%
\indexS{anguli alterni}
ergo ut \textit{EI} ad \textit{GE},
seu ut \textit{A} ad \textit{B},
ita \textit{AG} ad \textit{AN}.
Eodem modo~\textit{A}~ad~\textit{C}
ut \textit{AG} ad \textit{AK}.%
}%
}{%
\lemma{Si angulus \bracks{...} \textit{AG} ad \textit{AK}}%
\Cfootnote{%
Vgl. a.a.O., prop.\;XI %
(S.\;495b: Zitat mit Auslassungen).\cite{00124}%
}}
\edtext{%
Jam triangula \textit{EIG}, \textit{FHE} sunt similia,%
\indexS{triangula similia}
ergo \textit{FHE} %
ipsi \textit{GNA},
ergo $HE : FE : FH\; \squaredots\; AN : AG : NG$
seu $AK \squaredots B.\:A.\:C$.%
}{%
\lemma{Jam triangula \bracks{...} \textit{B.\:A.\:C}}%
\Cfootnote{%
Vgl. \textsc{Dechales}, \textit{Cursus}, tract.\:VII, lib.\;III, prop.\;XI, coroll. %
(S.\;495b).\cite{00124}%
}}
\pend%
\begin{edfigure}[1.25em][1.00em]%
\hfill%
\includegraphics[scale=0.90]{%
gesamttex/edit_VIII-4/images/LH_35_09_17_005-006_d16_006r.pdf}%
\hfill%
\includegraphics[scale=0.90]{%
gesamttex/edit_VIII-4/images/LH_35_09_17_005-006_d17_006r.pdf}%
\hfill%
\edcaption[0.25em]{%
\hspace*{-5mm}%
\lbrack\textit{Fig.\:5}\rbrack%
\edtext{}{%
\lemma{\lbrack\textit{Fig.\:5}\rbrack}%
\Cfootnote[nonum]{%
Vgl. a.a.O., prop.\;XI %
(S.\;495b: einziges Diagramm).\cite{00124}%
}}%
\hspace*{50mm}
\lbrack\textit{Fig.\:6}\rbrack%
\edtext{}{%
\lemma{\lbrack\textit{Fig.\:6}\rbrack}%
\Cfootnote[nonum]{%
Vgl. a.a.O., prop.\;XII %
(S.\;496a: einziges Diagramm).\cite{00124}%
}}%
}
\end{edfigure}%
\pstart%
\edtext{%
Sit
\edtext{%
angulus \textit{A} obtusus.%
\indexS{angulus obtusus}%
}{%
\lemma{angulus \textit{A}}%
\Bfootnote{%
\textit{(1)}~acutus %
\textit{(2)}~obtusus%
~\textit{L}%
}%
\lemma{angulus \textit{A} obtusus}%
\Cfootnote{%
Vgl. das Diagramm \bracks{\textit{Fig.\:6}}.%
}%
}
Caetera ut ante.
$C:A\squaredots FG:FI$.
Jam $\triangledown$\textsuperscript{la} \textit{FIE} simil.\;\textit{FGH}.%
\indexS{triangula similia}
$FG:FI \squaredots FH:FE$.
Ergo haec\:$\squaredots\: C:A$.
Eodem modo $B:A \squaredots HE:FE$.
Ergo $B:A:C$\;$\squaredots$ $HE:FE:HF$.
Ergo et
(ob sim. $\triangledown$ \textit{FIE}, \textit{FGH})
$ \squaredots\; FL : HE : IE$.
Denique et triangula similia sunt \textit{FHE}, \textit{GNA},%
\indexS{triangula similia}
\textit{angulus enim \textit{GNA}
duobus internis \textit{NFG}, \textit{NGF} seu \textit{AEG},
hoc est angulo \textit{FHE},
aequalis est.
Item angulus \textit{NAG}
duobus internis \textit{ALG} seu \textit{AEG}
et \textit{AGL} seu \textit{AEL},
hoc est toti \textit{FEH} aequalis est.
Eodem modo ostendam
triangulum \textit{GAK} eidem}
\edtext{%
\textit{FHE}
aequaleangulum.%
\indexS{triangulum aequaleangulum}%
}{%
\lemma{\textit{FHE}}%
\Bfootnote{%
\textit{(1)}~aequi %
\textit{(2)}~aequaleangulum.~%
~\textit{L}%
}}
Ergo $HE:FE:HF \squaredots NA:AG:NG$ seu \textit{AK}.%
}{%
\lemma{Sit angulus \bracks{...} seu \textit{AK}}%
\Cfootnote{%
Vgl. a.a.O. %
(S.\;496a: Zitat mit Auslassungen).\cite{00124}%
}}
\edtext{%
Ex his converse
\edtext{%
determ\bracks{inab}iles angul\bracks{i}%
}{%
\lemma{determaniles angulos}%
\Bfootnote{%
\textit{L\;ändert Hrsg.}%
}}
datis ponderibus.%
}{%
\lemma{Ex his \bracks{...} ponderibus}%
\Cfootnote{%
Vgl. \textsc{Dechales}, \textit{Cursus}, tract.\:VII, lib.\;III, prop.\;XII, coroll. %
(S.\;496a).\cite{00124}%
}}
\pend%
\begin{edfigure}[0.00em][1.75em]%
\hfill%
\includegraphics[scale=0.95]{%
gesamttex/edit_VIII-4/images/LH_35_09_17_005-006_d18_006r.pdf}
\hfill%
\includegraphics[scale=0.95]{%
gesamttex/edit_VIII-4/images/LH_35_09_17_005-006_d19_006r.pdf}
\hfill%
\edcaption[0.25em]{%
\lbrack\textit{Fig.\:7}\rbrack%
\hspace*{35mm}%
\lbrack\textit{Fig.\:8}\rbrack%
\edtext{}{%
\lemma{\bracks{\textit{Fig.\:7}}}%
\Cfootnote[nonum]{%
Vgl. a.a.O., prop.\;XIII %
(S.\;496b: erstes Diagramm).\cite{00124}%
}}%
\edtext{}{%
\lemma{\bracks{\textit{Fig.\:8}}}%
\Cfootnote[nonum]{%
Vgl. a.a.O. %
(S.\;496b: zweites Diagramm).\cite{00124}%
}}%
}
\end{edfigure}%
\pstart%
\edtext{%
\textit{Si
\edtext{%
grave
\edtext{\textit{AB}%
}{%
\lemma{\textit{AB}}%
\Bfootnote{%
\textit{erg.\;L}%
}}%
}{%
\lemma{\textit{grave AB}}%
\Cfootnote{%
Vgl. das Diagramm \bracks{\textit{Fig.\:7}}.%
}}
e duobus}
\edtext{%
\textit{funibus}%
\indexS{grave pendens e duobus funibus}%
\indexS{funis tensus}
\bracks{\textit{FE}}%
\textit{, \textit{DC}
non}%
}{%
\lemma{\textit{funibus}}%
\Bfootnote{%
\bars{%
\bars{\textit{FB} \textit{ändert Hrsg. nach Vorlage}}%
\textit{, DC}
\textit{erg.}} \textit{non}%
~\textit{L}%
}}
\textit{parallelis pendeat},
et \textit{ex} eorum \textit{concursu}\;\textit{G}%
\indexS{concursus funium}
\textit{linea perpendiculi}%
\indexS{linea perpendiculi}%
\indexS{perpendiculum}
\edtext{%
\textit{GH} \textit{ducatur},%
}{%
\lemma{\textit{GH}}%
\Bfootnote{%
\textit{(1)}~demittatur %
\textit{(2)}~\textit{ducatur}%
~\textit{L}%
}}
ea \textit{transibit per centrum gravitatis}%
\indexS{centrum gravitatis gravis pendentis}
ipsius gravis \textit{AB}.
Nam grave eodem modo sustentatur ex puncto \textit{C},%
\indexS{grave sustentatum}
quo ex \textit{G},
et ex puncto
\edtext{%
\lbrack\textit{E}\rbrack%
}{%
\lemma{\textit{e}}%
\Bfootnote{%
\textit{L\;ändert Hrsg.}%
}}
quo ex eodem\;\textit{G}.
Jam si ex \textit{G} suspensum esset,%
\indexS{grave suspensum}
linea perpendiculi%
\indexS{linea perpendiculi}%
\indexS{perpendiculum}
transiret per centrum gravitatis.%
\indexS{centrum gravitatis gravis suspensi}%
}{%
\lemma{\textit{Si grave} \bracks{...} gravitatis}%
\Cfootnote{%
Vgl. a.a.O. %
(S.\;496b: Zitat mit Auslassungen).\cite{00124}%
}}
\pend%
\pstart%
(+\;Hoc etiam sic demonstres:
Eodem modo sustentatur ex
\edtext{%
\lbrack\textit{CE}\rbrack,%
}{%
\lemma{\textit{ce}}%
\Bfootnote{%
\textit{L\;ändert Hrsg.}%
}}
quomodo ex \textit{LM},
posito \textit{LM} ipsi \textit{CE} parallela,
et omnibus similibus
et similiter
\edtext{%
positis,
et si \textit{GM} vel \textit{GL} ponatur
infinite parva,%
}{%
\lemma{positis,}%
\Bfootnote{%
\textit{(1)}~jam si \textit{LM}, ponatur in %
\textit{(2)}~et si %
\bracks{...} infinite parva,%
~\textit{L}%
}}
nihilo minus eodem modo sustinebitur%
\indexS{grave sustentum}%
\bracks{;}
ut eo
\edtext{%
casu
(quo scilicet revera
suspenditur \textit{AB}
ex ipso \textit{G})
transit%
}{%
\lemma{casu}%
\Bfootnote{%
\textit{(1)}~transit %
\textit{(2)}~(quo scilicet %
\bars{revera \textit{erg.}} suspenditur \textit{AB} ex ipso \textit{G}) transit%
~\textit{L}%
}}
linea directionis%
\indexS{linea directionis}
per centrum gravitatis,%
\indexS{centrum gravitatis gravis suspensi}
ergo et alias.%
\;\bracks{+)}
\pend%
\selectlanguage{latin}%
\frenchspacing%
\normalsize%
